\documentclass[10pt]{extarticle}
\usepackage[a4paper,margin=1.5cm]{geometry}
\usepackage[T1]{fontenc}
\usepackage{amsmath,amssymb}
\usepackage{graphicx}
\usepackage[export]{adjustbox}
\usepackage{multicol}
\usepackage{enumitem}
\usepackage{booktabs,array,tabularx}
\setlength{\columnsep}{0.7cm}
\setlength{\parindent}{0pt}
\setlength{\parskip}{3pt}
\pagestyle{empty}
\makeatletter
% figure registry (figdefs.tex is generated by the extract step)
\newcommand{\qfigdef}[4]{\@namedef{qf@file@#1}{#2}\@namedef{qf@w@#1}{#3}\@namedef{qf@wide@#1}{#4}}
\newcommand{\qf@check}[1]{\@ifundefined{qf@file@#1}{\PackageError{score2}{Unknown figure id #1}{}}{}}
\newcommand{\qf@img}[4][]{\includegraphics[width=#2\dimexpr\@nameuse{qf@w@#4}pt\relax,max width=#3,max height=0.45\textheight,#1]{\@nameuse{qf@file@#4}}}
% Full-width questions: finalize wraps every question that contains a wide figure in \qfullbegin...\qfullend,
% so the WHOLE question (stem, figures, options) is set across the page in reading order, then the two columns
% resume. The question's head - its number, stem and first figure (\qfullsplit marks the end of the head) -
% is kept on one page, so a number is never left at a page foot with its figure overleaf; the rest may break
% across pages normally, so no big gaps are left. Without a \qfullsplit, the whole question is kept together
% if it is under 45% of a page. Never write these yourself.
\newif\ifqf@full
\newif\ifqf@boxed
\newsavebox{\qf@box}
\newcommand{\qf@open}{\global\qf@fulltrue\par\global\qf@boxedtrue
  \setbox\qf@box\vbox\bgroup\hsize\textwidth\linewidth\textwidth}
\newcommand{\qf@emit}[1]{\par\egroup\global\qf@boxedfalse
  \ifdim\dimexpr\ht\qf@box+\dp\qf@box\relax<#1\textheight\noindent\usebox{\qf@box}\par\else\unvbox\qf@box\fi}
\newcommand{\qfullsplit}{\ifqf@boxed\qf@emit{0.85}\fi}
\newcommand{\qf@close}{\ifqf@boxed\qf@emit{0.45}\fi\global\qf@fullfalse}
\newcommand{\qfullbegin}{\end{multicols}\qf@open}
\newcommand{\qfullend}{\qf@close\begin{multicols}{2}}
\newcommand{\qfullstart}{\qf@open}
\newcommand{\qfulldone}{\qf@close}
% \qfig[scale]{id}: block figure, centred, directly where it is written (size set by the tool).
% Fallback if a wide figure is not inside a full-width question: break out of the columns in place - no floats.
\newcommand{\qfig}[2][1]{\qf@check{#2}%
  \ifqf@full\qf@blk{#1}{#2}\else
  \ifnum\@nameuse{qf@wide@#2}=1 \expandafter\@firstoftwo\else\expandafter\@secondoftwo\fi
  {\end{multicols}\begin{center}\qf@img{#1}{\textwidth}{#2}\end{center}\begin{multicols}{2}}%
  {\qf@blk{#1}{#2}}\fi}
\newcommand{\qf@blk}[2]{\par\smallskip{\centering\qf@img{#1}{\linewidth}{#2}\par}\smallskip}
% \qfigi[scale]{id}: inline figure (inside option lists / tables / minipages); never breaks columns.
\newcommand{\qfigi}[2][1]{\qf@check{#2}\qf@img[valign=c]{#1}{\linewidth}{#2}}
\makeatother
% structure macros - lint reads these to check completeness and figure placement
\newcommand{\papersec}[1]{\par\bigskip{\centering\Large\bfseries PAPER-#1\par}\medskip}
\newcommand{\partsec}[2]{\par\medskip\noindent{\large\bfseries PART-#1: #2}\par\smallskip}
\newcommand{\sect}[1]{\par\smallskip\noindent\textbf{#1}\par}
\newcommand{\qn}[1]{\par\smallskip\noindent\textbf{#1.}~}
\newcommand{\paraq}[2]{\par\smallskip\noindent\textbf{Paragraph for Questions #1 and #2}\par}
\newenvironment{qoptions}{\begin{itemize}[leftmargin=1.8em,itemsep=1pt,topsep=2pt,label={}]}{\end{itemize}}
\newcommand{\opt}[2]{\item[(#1)] #2}
% figure registry (generated by score2_tools.py finalize)
\qfigdef{p1_phy_s1_q01_1}{Internal Test - Test 03 (08 Sep 2021)-img/p1_phy_s1_q01_1.png}{278.8}{1}
\qfigdef{p1_phy_s1_q02_1}{Internal Test - Test 03 (08 Sep 2021)-img/p1_phy_s1_q02_1.png}{94.5}{0}
\qfigdef{p1_phy_s1_q05_1}{Internal Test - Test 03 (08 Sep 2021)-img/p1_phy_s1_q05_1.png}{111.3}{0}
\qfigdef{p1_phy_s1_q06_1}{Internal Test - Test 03 (08 Sep 2021)-img/p1_phy_s1_q06_1.png}{126.1}{0}
\qfigdef{p1_phy_s1_q07_1}{Internal Test - Test 03 (08 Sep 2021)-img/p1_phy_s1_q07_1.png}{38.6}{0}
\qfigdef{p1_phy_s1_q07_2}{Internal Test - Test 03 (08 Sep 2021)-img/p1_phy_s1_q07_2.png}{38.6}{0}
\qfigdef{p1_phy_s1_q08_1}{Internal Test - Test 03 (08 Sep 2021)-img/p1_phy_s1_q08_1.png}{79.8}{0}
\qfigdef{p1_phy_s1_q08_2}{Internal Test - Test 03 (08 Sep 2021)-img/p1_phy_s1_q08_2.png}{75.7}{0}
\qfigdef{p1_phy_s1_q08_3}{Internal Test - Test 03 (08 Sep 2021)-img/p1_phy_s1_q08_3.png}{80.3}{0}
\qfigdef{p1_phy_s1_q08_4}{Internal Test - Test 03 (08 Sep 2021)-img/p1_phy_s1_q08_4.png}{80.5}{0}
\qfigdef{p1_phy_s1_q09_1}{Internal Test - Test 03 (08 Sep 2021)-img/p1_phy_s1_q09_1.png}{134.8}{0}
\qfigdef{p1_phy_s1_q10_1}{Internal Test - Test 03 (08 Sep 2021)-img/p1_phy_s1_q10_1.png}{41.6}{0}
\qfigdef{p1_phy_s1_q11_1}{Internal Test - Test 03 (08 Sep 2021)-img/p1_phy_s1_q11_1.png}{88.5}{0}
\qfigdef{p1_phy_s2_q13_1}{Internal Test - Test 03 (08 Sep 2021)-img/p1_phy_s2_q13_1.png}{89.3}{0}
\qfigdef{p1_phy_s2_q14_1}{Internal Test - Test 03 (08 Sep 2021)-img/p1_phy_s2_q14_1.png}{208.5}{0}
\qfigdef{p1_phy_s2_q16_1}{Internal Test - Test 03 (08 Sep 2021)-img/p1_phy_s2_q16_1.png}{207.3}{0}
\qfigdef{p1_phy_s2_q16_2}{Internal Test - Test 03 (08 Sep 2021)-img/p1_phy_s2_q16_2.png}{163.2}{0}
\qfigdef{p1_chem_s1_q01_1}{Internal Test - Test 03 (08 Sep 2021)-img/p1_chem_s1_q01_1.png}{181.3}{0}
\qfigdef{p1_chem_s1_q01_A}{Internal Test - Test 03 (08 Sep 2021)-img/p1_chem_s1_q01_A.png}{99.3}{0}
\qfigdef{p1_chem_s1_q01_B}{Internal Test - Test 03 (08 Sep 2021)-img/p1_chem_s1_q01_B.png}{167.1}{0}
\qfigdef{p1_chem_s1_q01_C}{Internal Test - Test 03 (08 Sep 2021)-img/p1_chem_s1_q01_C.png}{89.9}{0}
\qfigdef{p1_chem_s1_q01_D}{Internal Test - Test 03 (08 Sep 2021)-img/p1_chem_s1_q01_D.png}{91.5}{0}
\qfigdef{p1_chem_s1_q02_1}{Internal Test - Test 03 (08 Sep 2021)-img/p1_chem_s1_q02_1.png}{180.0}{0}
\qfigdef{p1_chem_s1_q02_A}{Internal Test - Test 03 (08 Sep 2021)-img/p1_chem_s1_q02_A.png}{130.9}{0}
\qfigdef{p1_chem_s1_q02_B}{Internal Test - Test 03 (08 Sep 2021)-img/p1_chem_s1_q02_B.png}{129.6}{0}
\qfigdef{p1_chem_s1_q02_C}{Internal Test - Test 03 (08 Sep 2021)-img/p1_chem_s1_q02_C.png}{86.8}{0}
\qfigdef{p1_chem_s1_q02_D}{Internal Test - Test 03 (08 Sep 2021)-img/p1_chem_s1_q02_D.png}{84.6}{0}
\qfigdef{p1_chem_s1_q03_1}{Internal Test - Test 03 (08 Sep 2021)-img/p1_chem_s1_q03_1.png}{217.3}{0}
\qfigdef{p1_chem_s2_q05_A}{Internal Test - Test 03 (08 Sep 2021)-img/p1_chem_s2_q05_A.png}{72.5}{0}
\qfigdef{p1_chem_s2_q05_B}{Internal Test - Test 03 (08 Sep 2021)-img/p1_chem_s2_q05_B.png}{90.1}{0}
\qfigdef{p1_chem_s2_q05_D}{Internal Test - Test 03 (08 Sep 2021)-img/p1_chem_s2_q05_D.png}{84.5}{0}
\qfigdef{p1_chem_s2_q07_A}{Internal Test - Test 03 (08 Sep 2021)-img/p1_chem_s2_q07_A.png}{210.5}{0}
\qfigdef{p1_chem_s2_q07_2}{Internal Test - Test 03 (08 Sep 2021)-img/p1_chem_s2_q07_2.png}{148.4}{0}
\qfigdef{p1_chem_s2_q07_C}{Internal Test - Test 03 (08 Sep 2021)-img/p1_chem_s2_q07_C.png}{140.3}{0}
\qfigdef{p1_chem_s2_q07_D}{Internal Test - Test 03 (08 Sep 2021)-img/p1_chem_s2_q07_D.png}{163.6}{0}
\qfigdef{p1_chem_s2_q08_A}{Internal Test - Test 03 (08 Sep 2021)-img/p1_chem_s2_q08_A.png}{58.7}{0}
\qfigdef{p1_chem_s2_q08_B}{Internal Test - Test 03 (08 Sep 2021)-img/p1_chem_s2_q08_B.png}{78.5}{0}
\qfigdef{p1_chem_s2_q08_C}{Internal Test - Test 03 (08 Sep 2021)-img/p1_chem_s2_q08_C.png}{61.3}{0}
\qfigdef{p1_chem_s2_q08_D}{Internal Test - Test 03 (08 Sep 2021)-img/p1_chem_s2_q08_D.png}{57.2}{0}
\qfigdef{p1_chem_s2_q09_1}{Internal Test - Test 03 (08 Sep 2021)-img/p1_chem_s2_q09_1.png}{282.1}{1}
\qfigdef{p1_chem_s2_q09_A}{Internal Test - Test 03 (08 Sep 2021)-img/p1_chem_s2_q09_A.png}{46.0}{0}
\qfigdef{p1_chem_s2_q09_2}{Internal Test - Test 03 (08 Sep 2021)-img/p1_chem_s2_q09_2.png}{80.1}{0}
\qfigdef{p1_chem_s2_q09_C}{Internal Test - Test 03 (08 Sep 2021)-img/p1_chem_s2_q09_C.png}{61.5}{0}
\qfigdef{p1_chem_s2_q09_4}{Internal Test - Test 03 (08 Sep 2021)-img/p1_chem_s2_q09_4.png}{26.8}{0}
\qfigdef{p1_chem_s2_q10_1}{Internal Test - Test 03 (08 Sep 2021)-img/p1_chem_s2_q10_1.png}{128.9}{0}
\qfigdef{p1_chem_s2_q11_1}{Internal Test - Test 03 (08 Sep 2021)-img/p1_chem_s2_q11_1.png}{124.8}{0}
\qfigdef{p1_chem_s2_q11_A}{Internal Test - Test 03 (08 Sep 2021)-img/p1_chem_s2_q11_A.png}{70.9}{0}
\qfigdef{p1_chem_s2_q11_B}{Internal Test - Test 03 (08 Sep 2021)-img/p1_chem_s2_q11_B.png}{84.9}{0}
\qfigdef{p1_chem_s2_q11_C}{Internal Test - Test 03 (08 Sep 2021)-img/p1_chem_s2_q11_C.png}{52.4}{0}
\qfigdef{p1_chem_s2_q11_D}{Internal Test - Test 03 (08 Sep 2021)-img/p1_chem_s2_q11_D.png}{72.9}{0}
\qfigdef{p1_chem_s2_q12_A}{Internal Test - Test 03 (08 Sep 2021)-img/p1_chem_s2_q12_A.png}{286.5}{1}
\qfigdef{p1_chem_s2_q12_B}{Internal Test - Test 03 (08 Sep 2021)-img/p1_chem_s2_q12_B.png}{245.0}{0}
\qfigdef{p1_chem_s2_q12_C}{Internal Test - Test 03 (08 Sep 2021)-img/p1_chem_s2_q12_C.png}{259.8}{1}
\qfigdef{p1_chem_s2_q12_1}{Internal Test - Test 03 (08 Sep 2021)-img/p1_chem_s2_q12_1.png}{276.2}{1}
\qfigdef{p1_chem_s3_q13_1}{Internal Test - Test 03 (08 Sep 2021)-img/p1_chem_s3_q13_1.png}{201.2}{0}
\qfigdef{p1_chem_s3_q14_1}{Internal Test - Test 03 (08 Sep 2021)-img/p1_chem_s3_q14_1.png}{198.5}{0}
\qfigdef{p1_chem_s3_q15_1}{Internal Test - Test 03 (08 Sep 2021)-img/p1_chem_s3_q15_1.png}{169.1}{0}
\qfigdef{p1_chem_s3_q16_1}{Internal Test - Test 03 (08 Sep 2021)-img/p1_chem_s3_q16_1.png}{117.4}{0}
\qfigdef{p1_chem_s3_q18_1}{Internal Test - Test 03 (08 Sep 2021)-img/p1_chem_s3_q18_1.png}{189.6}{0}
\raggedcolumns
\makeatletter\@beginparpenalty=10000\makeatother
\newcolumntype{L}[1]{>{\raggedright\arraybackslash}p{#1}}
\begin{document}
\begin{center}{\Large\bfseries Internal Test --- Test 03 (08 Sep 2021)}\end{center}
\begin{multicols}{2}
\raggedcolumns
\qfullbegin
\papersec{1}
\partsec{1}{PHYSICS}
\sect{SECTION-I(i)}
\qn{1} The switch in circuit shifts from 1 to 2 when $V_C>2V/3$ and goes back to 1 from 2 when $V_C<V/3$. The voltmeter reads voltage as plotted. What is the period $T$ of the wave form in terms of $R$ and $C$?\par
\qfig{p1_phy_s1_q01_1}
\qfullsplit
\begin{qoptions}
\opt{A}{$RC\ln 3$}
\opt{B}{$2RC\ln 2$}
\opt{C}{$\dfrac{RC}{2}\ln 3$}
\opt{D}{$\dfrac{RC}{3}\ln 3$}
\end{qoptions}
\qfullend
\qn{2} A conducting rod of length $l$ and mass $m$ is moving down a smooth inclined plane of inclination $\theta$ with constant velocity $v$ in fig. A current $I$ is flowing in the conductor in a direction perpendicular to paper inwards. A vertically upward magnetic field $\vec B$ exists in space. Then magnitude of magnetic field $\vec B$ is\par
\qfig{p1_phy_s1_q02_1}
\begin{qoptions}
\opt{A}{$\dfrac{mg}{il}\sin\theta$}
\opt{B}{$\dfrac{mg}{il}\tan\theta$}
\opt{C}{$\dfrac{mg\cos\theta}{il}$}
\opt{D}{$\dfrac{mg}{il\sin\theta}$}
\end{qoptions}
\qn{3} If any values of current gives a $45^\circ$ deflection in property arranged tangent galvanometer. If current becomes $\dfrac1{\sqrt3}$ of its initial value, then deflection in it :-
\begin{qoptions}
\opt{A}{Decrease by $30^\circ$}
\opt{B}{Decrease by $15^\circ$}
\opt{C}{Increase by $15^\circ$}
\opt{D}{Increase by $30^\circ$}
\end{qoptions}
\qn{4} In a 20 m deep lake, the bottom is at a constant temperature of $4^\circ$C. The air temperature is constant at $-10^\circ$C. The thermal conductivity of ice is 4 times that water. Neglecting the expansion of water on freezing, the maximum thickness of ice will be
\begin{qoptions}
\opt{A}{$\dfrac{20}{11}$ m}
\opt{B}{$\dfrac{200}{11}$ m}
\opt{C}{$20$ m}
\opt{D}{$10$ m}
\end{qoptions}
\sect{SECTION-I(ii)}
\qn{5} In the given figure, a capacitor of non-parallel plates is shown. The plates of capacitor are connected by a cell of emf $V_0$. If $\sigma$ denotes surface charge density and $E$ denotes electric field. Then\par
\qfig{p1_phy_s1_q05_1}
\begin{qoptions}
\opt{A}{$\sigma_A>\sigma_B$}
\opt{B}{$E_F>E_D$}
\opt{C}{$E_F=E_D$}
\opt{D}{$\sigma_A=\sigma_B$}
\end{qoptions}
\qn{6} Two capacitors of $2\,\mu$F and $3\,\mu$F are charged to 150 volt and 120 volt respectively. The plates of capacitor are connected as shown in the figure. A discharged capacitor of capacity $1.5\,\mu$F falls to the free ends of the wire. Then --\par
\qfig{p1_phy_s1_q06_1}
\begin{qoptions}
\opt{A}{Charge on the $1.5\,\mu$F capacitor is $180\,\mu$C}
\opt{B}{Charge on the $2\,\mu$F capacitor is $120\,\mu$F}
\opt{C}{Charge flows through A from right to left}
\opt{D}{Charge flows through A from left to right.}
\end{qoptions}
\qn{7} An observer A and a charge $Q$ are fixed in a stationary frame $F_1$. An observer B is fixed in a frame $F_2$, which is moving with respect to $F_1$.\par
\qfigi{p1_phy_s1_q07_1} \qfigi{p1_phy_s1_q07_2}
\begin{qoptions}
\opt{A}{Both A and B will observe electric fields.}
\opt{B}{Both A and B will observe magnetic fields.}
\opt{C}{Neither A nor B will observe magnetic fields.}
\opt{D}{B will observe a magnetic field, but A will not.}
\end{qoptions}
\qn{8} In the loops shown, all curved sections are either semicircles or quarter circles. All the loops carry the same current. The magnetic fields at the centres have magnitudes $B_1$, $B_2$, $B_3$ and $B_4$\par
\qfigi{p1_phy_s1_q08_1} \qfigi{p1_phy_s1_q08_2} \qfigi{p1_phy_s1_q08_3} \qfigi{p1_phy_s1_q08_4}
\begin{qoptions}
\opt{A}{$B_4$ is maximum.}
\opt{B}{$B_3$ is minimum.}
\opt{C}{$B_4>B_1>B_2>B_3$}
\opt{D}{$B_1>B_4>B_3>B_2$}
\end{qoptions}
\qn{9} A particle of charge $+q$ and mass $m$ moving under the influence of a uniform electric field $E\hat i$ and uniform magnetic field $B\hat k$ follows a trajectory from P to Q as shown in figure. The velocities at P and Q are $v\hat i$ and $-2v\hat j$ respectively. Which of the following statement(s) is/are correct ?\par
\qfig{p1_phy_s1_q09_1}
\begin{qoptions}
\opt{A}{$E=\dfrac34\left[\dfrac{mv^2}{qa}\right]$}
\opt{B}{Rate of work done by the electric field at P is $\dfrac34\left[\dfrac{mv^3}{a}\right]$}
\opt{C}{Rate of work done by the electric field at P is zero}
\opt{D}{Rate of work done by both the fields at Q is zero}
\end{qoptions}
\qn{10} From a cylinder of radius $R$, a cylinder of radius $R/2$ is removed, as shown. Current flowing in the remaining cylinder is $I$. Magnetic field strength is--\par
\qfig{p1_phy_s1_q10_1}
\begin{qoptions}
\opt{A}{zero at point A}
\opt{B}{zero at point B}
\opt{C}{$\dfrac{\mu_0I}{3\pi R}$ at point A}
\opt{D}{$\dfrac{\mu_0I}{3\pi R}$ at point B}
\end{qoptions}
\qn{11} A triangular plate has two cavities, one square and one rectangular as shown in figure. The plate is heated.\par
\qfig{p1_phy_s1_q11_1}
\begin{qoptions}
\opt{A}{$a$ increase, $b$ decrease}
\opt{B}{$a$ and $b$ both increase}
\opt{C}{$a$ and $b$ increase, $x$ and $l$ decrease}
\opt{D}{$a$, $b$, $x$ and $l$ all increase}
\end{qoptions}
\qn{12} A hollow copper sphere \& a hollow copper cube of same surface area \& negligible thickness, are filled with warm water of same temperature and placed in an enclosure of constant temperature, a few degrees below that of the bodies. Then in the beginning :
\begin{qoptions}
\opt{A}{The rate of energy lost by the sphere is greater than that by the cube}
\opt{B}{The rate of energy lost by the two are equal}
\opt{C}{The rate of energy lost by the sphere is less than that by the cube}
\opt{D}{The rate of fall of temperature for sphere is less than that for the cube.}
\end{qoptions}
\sect{SECTION-II}
\qn{13} Two medium of dielectric constant $K_1$ and $K_2$ are introduced according to given figure. If $\dfrac{K_1}{K_2}=3$ then calculate ratio of capacity of part A and part B. (Area of each plate is $A$)\par
\qfig{p1_phy_s2_q13_1}
\qn{14} Two short magnet A and B of magnetic dipole moments $M_1$ and $M_2$ respectively are placed as shown. The axis of `A' and the equatorial line of `B' are the same. The magnetic force on one magnet due to the other is $\dfrac{-K\mu_0M_1M_2}{4\pi r^4}$ find $k$ ?\par
\qfig{p1_phy_s2_q14_1}
\qn{15} A steady current $I$ goes through a wire loop PQR having shape of a right angle triangle with $PQ=3x$, $PR=4x$ and $QR=5x$. If the magnitude of the magnetic field at P due to this loop is $k\left(\dfrac{\mu_0I}{48\pi x}\right)$, find the value of $k$.
\qn{16} Two different rods A and B are kept as shown in figure.\par
\qfig{p1_phy_s2_q16_1}
\qfig{p1_phy_s2_q16_2}
The variation of temperature of different cross sections is plotted in a graph shown in figure. Find the ratio of thermal conductivities of B to A.
\qn{17} A body cools from $50^\circ$C to $49.9^\circ$C in 5 sec. It cools from $40^\circ$C to $39.9^\circ$C in $t$ sec. Assuming Newtons law of cooling to be valid and temperature of surrounds at $30^\circ$C, value of $t/5$ will be? (in sec) (to the nearest integer)
\qn{18} A metal rod AB of length $10x$ has its one end A in ice at $0^\circ$C and the other end B in water at $100^\circ$C. It a point P on the rod is maintained at $400^\circ$C, then it is found that equal amounts of water and ice evaporate and melt per unit time respectively. The latent heat of evaporation of water 540 cal/g and latent heat of melting of ice is 80 cal/g. If the point P is at $lx$ distance from the ice end then find out the value of $l$. [Neglect any heat loss to the surrounding.]
\partsec{2}{CHEMISTRY}
\sect{SECTION-I(i)}
\qn{1} What is the major product of following reaction\par
\qfig{p1_chem_s1_q01_1}
\begin{qoptions}
\opt{A}{\qfigi{p1_chem_s1_q01_A}}
\opt{B}{\qfigi{p1_chem_s1_q01_B}}
\opt{C}{\qfigi{p1_chem_s1_q01_C}}
\opt{D}{\qfigi{p1_chem_s1_q01_D}}
\end{qoptions}
\qn{2} What are the products X and Y.\par
\qfig{p1_chem_s1_q02_1}
\begin{qoptions}
\opt{A}{\qfigi{p1_chem_s1_q02_A}}
\opt{B}{\qfigi{p1_chem_s1_q02_B}}
\opt{C}{\qfigi{p1_chem_s1_q02_C}}
\opt{D}{\qfigi{p1_chem_s1_q02_D}}
\end{qoptions}
\qn{3} Which statement is correct about the end product of the following reaction series,\par
\qfig{p1_chem_s1_q03_1}
\begin{qoptions}
\opt{A}{It is optically inactive hydroxy ketone}
\opt{B}{It is a resolvable hydroxy acid}
\opt{C}{It is nonresolvable aldehyde}
\opt{D}{It is an optically inactive hydroxy acid}
\end{qoptions}
\qn{4} Which of the following compounds is strongest acid?
\begin{qoptions}
\opt{A}{Carbolic acid}
\opt{B}{Picric acid}
\opt{C}{Salycilic acid}
\opt{D}{Formic acid}
\end{qoptions}
\sect{SECTION-I(ii)}
\qn{5} Which of the following compounds are able to give red precipitate with Fehling's solution.
\begin{qoptions}
\opt{A}{\qfigi{p1_chem_s2_q05_A}}
\opt{B}{\qfigi{p1_chem_s2_q05_B}}
\opt{C}{$\mathrm{H{-}CO{-}H}$}
\opt{D}{\qfigi{p1_chem_s2_q05_D}}
\end{qoptions}
\qn{6} How many of the following reagents can be used for the conversion of R--CN into R--CHO
\begin{qoptions}
\opt{A}{DIBAL--H}
\opt{B}{$\mathrm{LiAlH_4}$ at higher temperatures}
\opt{C}{$\mathrm{SnCl_2/HCl}$}
\opt{D}{$\mathrm{H_2/Ni}$}
\end{qoptions}
\qn{7} Which of the following reactions are correctly matched with the products?
\begin{qoptions}
\opt{A}{\qfigi{p1_chem_s2_q07_A}}
\opt{B}{\qfigi{p1_chem_s2_q07_2}}
\opt{C}{\qfigi{p1_chem_s2_q07_C}}
\opt{D}{\qfigi{p1_chem_s2_q07_D}}
\end{qoptions}
\qn{8} Which of the following ethers can not be prepared by using williamson's synthesis.
\begin{qoptions}
\opt{A}{\qfigi{p1_chem_s2_q08_A}}
\opt{B}{\qfigi{p1_chem_s2_q08_B}}
\opt{C}{\qfigi{p1_chem_s2_q08_C}}
\opt{D}{\qfigi{p1_chem_s2_q08_D}}
\end{qoptions}
\qfullbegin
\qn{9} Which of the following are correct for the reaction.\par
\qfig{p1_chem_s2_q09_1}
\qfullsplit
\begin{qoptions}
\opt{A}{$\mathrm{A}=$ \qfigi{p1_chem_s2_q09_A}}
\opt{B}{$\mathrm{B}=$ \qfigi{p1_chem_s2_q09_2}}
\opt{C}{$\mathrm{C}=$ \qfigi{p1_chem_s2_q09_C}}
\opt{D}{$\mathrm{C}=$ \qfigi{p1_chem_s2_q09_4}}
\end{qoptions}
\qfullend
\qn{10} Choose the correct statement/s about, below mentioned reaction.\par
\qfig{p1_chem_s2_q10_1}
\begin{qoptions}
\opt{A}{Total three different carbanions are possible for the reactant}
\opt{B}{Gives two same sized cyclised products}
\opt{C}{Gives seven membered cyclised product as a minor product}
\opt{D}{Reaction gives intermolecular aldol condensation majorly.}
\end{qoptions}
\qn{11} What are the possible products of the following reaction.\par
\qfig{p1_chem_s2_q11_1}
\begin{qoptions}
\opt{A}{\qfigi{p1_chem_s2_q11_A}}
\opt{B}{\qfigi{p1_chem_s2_q11_B}}
\opt{C}{\qfigi{p1_chem_s2_q11_C}}
\opt{D}{\qfigi{p1_chem_s2_q11_D}}
\end{qoptions}
\qfullbegin
\qn{12} Which of the following reactions are correctly matched, with the name
\begin{qoptions}
\opt{A}{\qfigi{p1_chem_s2_q12_A}}
\opt{B}{\qfigi{p1_chem_s2_q12_B}}
\opt{C}{\qfigi{p1_chem_s2_q12_C}}
\opt{D}{\qfigi{p1_chem_s2_q12_1}}
\end{qoptions}
\qfullend
\sect{SECTION-II}
\qn{13} How many cyclic acetals are possible from the following reaction [including stereoisomers]\par
\qfig{p1_chem_s3_q13_1}
\qn{14} How many moles of HCHO is consumed in the following reaction\par
\qfig{p1_chem_s3_q14_1}
\qn{15} \qfig{p1_chem_s3_q15_1}
What is the degree of unsaturation in the product \textcircled{A} ?
\qn{16} \qfig{p1_chem_s3_q16_1}
(How many `A' are possible)
\qn{17} How many isomeric products would be possible from the following reaction.\par
$\mathrm{CH_3{-}CO{-}CO{-}CH_2CH_3}$ $\xrightarrow[\text{2. H}_2\text{O}]{\text{1. LiAlH}_4\text{/Dry Ether excess}}$
\qn{18} \qfig{p1_chem_s3_q18_1}
How many moles of $\mathrm{CH_3MgBr}$ would be consumed per mole of the reactant in the above reaction.
\partsec{3}{MATHEMATICS}
\sect{SECTION-I(i)}
\qn{1} Let $f(x)$ be a cubic polynomial on $R$ which increases in the interval $(-\infty,0)\cup(1,\infty)$ and decreases in the interval $(0,1)$. If $f'(2)=6$ and $f(2)=2$, then $f(0)=$
\begin{qoptions}
\opt{A}{$-2$}
\opt{B}{$0$}
\opt{C}{$2$}
\opt{D}{$1/2$}
\end{qoptions}
\qn{2} If $f(x)=\displaystyle\int_{x-1}^{2x}\left(t^3-2t+4\right)dt$ then $f'(0)=$
\begin{qoptions}
\opt{A}{$3$}
\opt{B}{$-3$}
\opt{C}{$0$}
\opt{D}{$1$}
\end{qoptions}
\qn{3} $\displaystyle\int\sec^{4/9}\theta\cdot\mathrm{cosec}^{14/9}\theta\,d\theta$ is equal to
\begin{qoptions}
\opt{A}{$\dfrac59(\tan\theta)^{-5/9}+C$}
\opt{B}{$-\dfrac95(\tan\theta)^{-5/9}+C$}
\opt{C}{$\dfrac95(\tan\theta)^{-9/5}+C$}
\opt{D}{$-\dfrac59(\tan\theta)^{-9/5}+C$}
\end{qoptions}
\qn{4} The value of $\displaystyle\int_0^{100}\left[\tan^{-1}x\right]dx$ is (where $[\,.\,]$ denotes greatest integer function)
\begin{qoptions}
\opt{A}{$100$}
\opt{B}{$100-\dfrac\pi4$}
\opt{C}{$100-\tan 1$}
\opt{D}{$100-\tan^{-1}\left(\dfrac\pi4\right)$}
\end{qoptions}
\sect{SECTION-I(ii)}
\qn{5} If $I_n=\displaystyle\int_0^{\pi/4}\tan^n x\,dx$, then $(n\in N,\,n\ge2)$
\begin{qoptions}
\opt{A}{$I_n+I_{n-2}=\dfrac1{n-1}$}
\opt{B}{$I_{n+1}+I_{n-1}=\dfrac1n$}
\opt{C}{$\dfrac1{n+1}<2I_n<\dfrac1{n-1}$}
\opt{D}{$\left(I_n+I_{n+2}\right)=\dfrac1{n+1}$}
\end{qoptions}
\qn{6} Which of the following is/are true?
\begin{qoptions}
\opt{A}{$\sqrt{e^\pi}>\left(\dfrac\pi2\right)^e$}
\opt{B}{$e^\pi>\pi^e$}
\opt{C}{$e^3>3^e$}
\opt{D}{$e^2>2^e$}
\end{qoptions}
\qn{7} Let $\displaystyle\int\dfrac{x^{1/2}}{\sqrt{1-x^3}}dx=\dfrac23\,g\circ f(x)+c$ then
\begin{qoptions}
\opt{A}{$f(x)=\sqrt x$}
\opt{B}{$f(x)=x^{3/2}$}
\opt{C}{$f(x)=x^{2/3}$}
\opt{D}{$g(x)=\sin^{-1}x$}
\end{qoptions}
\qn{8} If $\displaystyle\int\dfrac{xe^x}{\sqrt{1+e^x}}dx=f(x)\cdot\sqrt{1+e^x}-2\ln\left(g(x)\right)+c$ then
\begin{qoptions}
\opt{A}{$f(x)=x-1$}
\opt{B}{$g(x)=\dfrac{\sqrt{1+e^x}-1}{\sqrt{1+e^x}+1}$}
\opt{C}{$g(x)=\dfrac{\sqrt{1+e^x}+1}{\sqrt{1+e^x}-1}$}
\opt{D}{$f(x)=2(x-2)$}
\end{qoptions}
\qn{9} If for $x\ge0$, $f(x)=\displaystyle\int_0^x\left(\dfrac{4^t-1}{2^t-1}\right)dt$, then which of the following is/are NOT true?
\begin{qoptions}
\opt{A}{$f(x)$ has a maxima at $x=0$}
\opt{B}{$f(x)$ has a minima at $x=0$}
\opt{C}{$f(x)$ has a maximum value $0$}
\opt{D}{$f(x)$ has minimum value $0$}
\end{qoptions}
\qn{10} Consider the function $f(x)=\begin{cases}\displaystyle\int_0^x\left(4+|t-2|\right)dt, & x>3\\ ax^2+bx, & x\le3\end{cases}$. If $f(x)$ is differentiable at $x=3$ then
\begin{qoptions}
\opt{A}{$a+b=\dfrac{83}{18}$}
\opt{B}{$a+b=\dfrac{85}{18}$}
\opt{C}{$ab=\dfrac7{27}$}
\opt{D}{$\dfrac ba=84$}
\end{qoptions}
\qn{11} The value of the integral $\int e^{\sin^2x}$ $\left(\cos x+\cos^3x\right)$ $\sin x\,dx$ is
\begin{qoptions}
\opt{A}{$\dfrac12e^{\sin^2x}\left(3-\sin^2x\right)+c$}
\opt{B}{$e^{\sin^2x}\left(1+\dfrac12\cos^2x\right)+c$}
\opt{C}{$e^{\sin^2x}\left(2+\cos^2x\right)+c$}
\opt{D}{$e^{\sin^2x}\left(2+\sin^2x\right)+c$}
\end{qoptions}
\qn{12} If $x_1$ \& $x_2$ are positive numbers between 0 \& 1, then which of the following are true?
\begin{qoptions}
\opt{A}{$\sin\left(\dfrac{x_1+x_2}2\right)\le\dfrac12\left(\sin x_1+\sin x_2\right)$}
\opt{B}{$\tan\left(\dfrac{x_1+x_2}2\right)\le\dfrac12\left(\tan x_1+\tan x_2\right)$}
\opt{C}{$\ln\left(\dfrac{x_1+x_2}2\right)\le\dfrac12\left(\ln x_1+\ln x_2\right)$}
\opt{D}{$\left(\dfrac{x_1+x_2}2\right)^2\le\dfrac{x_1^2+x_2^2}2$}
\end{qoptions}
\sect{SECTION-II}
\qn{13} The difference between the greatest \& least value of $f(x)=\displaystyle\int_0^x\left(6t^2-24\right)dt$ on $[1,3]$ is \underline{\hspace{1.5cm}}
\qn{14} If $f'(x)=\phi(x)$ and $\phi'(x)=f(x)\ \forall x\in R$. Also $f(3)=5$, $f'(3)=4$ then $\left(f(10)\right)^2-\left(\phi(10)\right)^2=$ \underline{\hspace{1.5cm}}
\qn{15} If $\displaystyle\int_{-\sqrt2}^{\sqrt2}\dfrac{3x^7+3x^6-9x^5-8x^3-12x^2+x+1}{x^2+2}dx$ $=\dfrac{\sqrt2}{20}\left(5\pi+K^2\right)$, then $K=$ \underline{\hspace{1.5cm}}
\qn{16} Let $f(x)=x^3+x^2+2x-1$, the minimum value of $[x]$ that satisfy $f\left(f(x)\right)>f(2x+1)$ is (where $[.]$: G.I.F) \underline{\hspace{1.5cm}}
\qn{17} If $f(x)=\displaystyle\int\dfrac{\sqrt{\tan x}}{\sin x\cdot\cos x}dx$ and $f(0)=3$, then $f\left(\dfrac\pi4\right)=$ \underline{\hspace{1.5cm}}
\qn{18} If $I=\displaystyle\int\sin^{-1}\left(\dfrac{2x}{1+x^2}\right)dx=\sqrt n\,x\tan^{-1}x-\ln\left(1+x^2\right)+c$, then $n=$ \underline{\hspace{1.5cm}}
\end{multicols}
\end{document}
